\section{Method}
\label{sec:method}

We present an output-distribution KL framework for training contextual sparsity predictors in SwiGLU-based language models (Figure~\ref{fig:overview}).
The method replaces per-neuron binary classification (BCE) with an output-distribution alignment objective that couples all pruning decisions through the model's forward pass and achieves competitive performance without auxiliary error compensation networks at inference.

\subsection{Problem Formulation}
\label{sec:method:formulation}

Consider a pretrained transformer with $L$ layers, each containing a SwiGLU feed-forward network~\cite{shazeer2020glu}.
For layer $\ell$, the FFN computes:
\begin{align}
\mathbf{y}^{(\ell)} = \bigl(\sigma(\mathbf{x}^{(\ell)} \mathbf{W}_g^{(\ell)}) \odot \mathbf{x}^{(\ell)} \mathbf{W}_u^{(\ell)}\bigr) \mathbf{W}_d^{(\ell)},
\label{eq:swiglu}
\end{align}
where $\mathbf{x}^{(\ell)} \in \mathbb{R}^{d}$ is the hidden state after attention ($d{=}4096$), $\sigma$ denotes SiLU, $\mathbf{W}_g^{(\ell)}, \mathbf{W}_u^{(\ell)} \in \mathbb{R}^{d \times d_{\text{ff}}}$ and $\mathbf{W}_d^{(\ell)} \in \mathbb{R}^{d_{\text{ff}} \times d}$ are the gate, up-projection, and down-projection matrices, and $d_{\text{ff}}{=}14{,}336$.

Define the gate activation $\mathbf{g}^{(\ell)} = \sigma(\mathbf{x}^{(\ell)} \mathbf{W}_g^{(\ell)}) \in \mathbb{R}^{d_{\text{ff}}}$.
Each element $g_j^{(\ell)}$ determines the contribution of neuron $j$ to the output.
We introduce a binary mask $\mathbf{m}^{(\ell)} \in \{0,1\}^{d_{\text{ff}}}$ that zeros out pruned neurons:
\begin{align}
\mathbf{y}_{\text{sparse}}^{(\ell)} = \bigl(\mathbf{m}^{(\ell)} \odot \mathbf{g}^{(\ell)} \odot \mathbf{x}^{(\ell)} \mathbf{W}_u^{(\ell)}\bigr) \mathbf{W}_d^{(\ell)}.
\label{eq:sparse_ffn}
\end{align}
Our goal is to learn a mask predictor $f_\theta^{(\ell)}(\mathbf{x}^{(\ell)}) \to \mathbf{s}^{(\ell)} \in (0,1)^{d_{\text{ff}}}$ that produces per-neuron importance scores from $\mathbf{x}^{(\ell)}$.
At inference, the scores are fed to a global top-$k$ allocation mechanism~\cite{liu2025teal} that distributes the sparsity budget across layers non-uniformly, retaining the $(1{-}\rho) \times L \times d_{\text{ff}}$ highest-scoring neurons model-wide for a target sparsity ratio $\rho$.


\subsection{Predictor Architecture}
\label{sec:method:predictor}

Each layer $\ell$ has an independent two-layer MLP predictor:
\begin{align}
f_\theta^{(\ell)}(\mathbf{x}^{(\ell)}) = \text{GELU}(\mathbf{x}^{(\ell)} \mathbf{W}_1^{(\ell)}) \, \mathbf{W}_2^{(\ell)},
\label{eq:predictor}
\end{align}
where $\mathbf{W}_1^{(\ell)} \in \mathbb{R}^{d \times d_b}$ and $\mathbf{W}_2^{(\ell)} \in \mathbb{R}^{d_b \times d_{\text{ff}}}$ with bottleneck dimension $d_b{=}128$.
The predictor takes $\mathbf{x}^{(\ell)}$ (the post-attention hidden state) as input, since it already encodes the token's contextual dependencies from self-attention.
The output is $d_{\text{ff}}$-dimensional logits $\mathbf{s}^{(\ell)}$, from which the mask is derived.

Sweeps over $d_b \in \{64, 128, 256\}$ show diminishing returns beyond 128, with $d_b{=}256$ adding 2$\times$ parameters for $<$0.1pp PPL gain (Appendix~\ref{app:db_sensitivity}).
We use GELU rather than ReLU to avoid dead predictor neurons during the early exploration phase when Gumbel noise dominates.
Each predictor has $(4096 \times 128) + (128 \times 14336) = 2.37\text{M}$ parameters; with 32 independent predictors (one per layer), the total overhead is 75.9M, less than 1\% of the 8B base.
Per-layer independence avoids sequential latency at inference; the $K$-curve analysis (\S\ref{sec:perlayer}) confirms that cross-layer coupling during training provides modest benefit beyond the output-distribution signal.


\subsection{Gumbel-Sigmoid Relaxation}
\label{sec:method:gumbel}

The hard mask $m_j = \mathbf{1}[s_j > \text{threshold}]$ is non-differentiable, preventing gradient flow from the KL objective through the mask to the predictor parameters.
We apply Gumbel-Sigmoid relaxation~\cite{maddison2017concrete,louizos2018l0,jang2017gumbel} to produce differentiable soft masks:
{\small
\begin{align}
\tilde{m}_j^{(\ell)} \!=\! \sigma_\tau\!\left(\!\log \frac{s_j^{(\ell)}}{1 {-} s_j^{(\ell)}} + \varepsilon_1 {-} \varepsilon_2\right)\!,\; \varepsilon_{1,2} {\sim} \text{Gumbel}(0,1),
\label{eq:gumbel}
\end{align}}
where $\sigma_\tau(x) = \text{sigmoid}(x / \tau)$ and $\tau$ is a temperature parameter.
The Gumbel noise provides stochastic exploration of mask configurations during training, while the temperature controls sharpness: as $\tau \to 0$, $\tilde{m}_j$ approaches a binary decision.

We anneal $\tau$ linearly from 1.0 to 0.1 over training.
At high temperature, the soft mask permits gradient flow through all neurons for broad exploration.
As temperature decreases, the mask concentrates on confident decisions, approximating the hard allocation used at inference.

\paragraph{Comparison with Straight-Through Estimation.}
The straight-through estimator (STE)~\cite{bengio2013ste} applies a hard threshold in the forward pass and passes gradients through unchanged ($\partial m / \partial s = 1$).
In our setting, STE produces discontinuous gradient signals: small parameter changes can flip many neurons simultaneously, destabilizing sparsity enforcement (actual 20.8\% at 30\% target).
Gumbel relaxation eliminates this pathology and contributes a 25.2\% improvement (PPL 15.77 $\to$ 11.79 at 50\% sparsity), as we verify empirically in \S\ref{sec:experiments}.


\subsection{KL Divergence Training Objective}
\label{sec:method:kl}

Let $\mathbf{z}_{\text{dense}}$ denote the output logits from a full dense forward pass, and $\mathbf{z}_{\text{sparse}}$ the logits from a forward pass with predictor-generated soft masks $\{\tilde{\mathbf{m}}^{(\ell)}\}_\ell$ applied at every layer.
We minimize the forward KL divergence between the dense and sparse output distributions:
\begin{align}
\mathcal{L}_{\text{KL}} \!=\! D_{\text{KL}}(p_{\text{dense}} \| p_{\text{sparse}}) \!=\! \sum_{v} p_{\text{dense}}^{(v)} \log \frac{p_{\text{dense}}^{(v)}}{p_{\text{sparse}}^{(v)}},
\label{eq:kl_loss}
\end{align}
where $p_{\text{dense}} = \text{softmax}(\mathbf{z}_{\text{dense}})$ and $p_{\text{sparse}} = \text{softmax}(\mathbf{z}_{\text{sparse}})$, summed over the vocabulary $v$ and averaged over all token positions in the sequence.
The base model is frozen; only the predictor parameters receive gradients.

We use forward KL (mass-covering) rather than reverse KL, as the goal is to preserve the dense model's primary predictions rather than explore alternative modes.
In sparsity prediction, mode-covering is critical: false negatives (omitting a rarely but critically active neuron) cascade through downstream layers under global allocation, whereas false positives merely waste capacity~\cite{wu2025rethinking}.
Forward KL outperforms reverse KL by 7.9\% and JSD by 6.0\% at 50\% (Table~\ref{tab:ablation}).
No temperature scaling ($T{=}1$) is applied~\cite{hinton2015distilling}.


\subsection{Geometric Penalty Schedule}
\label{sec:method:lagrangian}

We enforce a model-level sparsity target $\rho$ via constrained optimization with $\bar{s} = \frac{1}{L} \sum_\ell \frac{1}{d_{\text{ff}}} \sum_j \tilde{m}_j^{(\ell)}$ denoting the average retained fraction:
\begin{align}
\min_\theta \; \mathcal{L}_{\text{KL}} \quad \text{s.t.} \quad \bar{s} = 1 - \rho.
\label{eq:constrained}
\end{align}

\begin{figure}[t]
\centering
\includegraphics[width=\columnwidth]{lagrangian_dynamics.pdf}
\vspace{-8pt}
\caption{Geometric penalty schedule dynamics for 30\% sparsity target ($\lambda_{\max}{=}5000$). Left axis (log): $\lambda$ rises via geometric doubling once sparsity undershoots the target, saturating at its cap. Right axis: actual sparsity converges from natural density (${\approx}$0.41) toward 0.28.}
\label{fig:lagrangian_dynamics}
\vspace{-8pt}
\end{figure}

We relax this via a geometric penalty schedule:
\begin{align}
\mathcal{L} = \mathcal{L}_{\text{KL}} + \lambda \!\cdot\! \text{clamp}(|\bar{s} {-} (1{-}\rho)| {-} \epsilon,\, 0)^2,
\label{eq:lagrangian}
\end{align}
where $\epsilon = 0.02$ is a deadzone margin that suppresses oscillation near the target (sensitivity in Appendix~\ref{app:lagrangian}). The penalty weight $\lambda$ is updated via geometric doubling (doubled when sparsity exceeds target, halved otherwise), clamped to $[\lambda_{\text{floor}}, \lambda_{\max}]$, reaching target sparsity within 40--110 steps across all tested regimes (Figure~\ref{fig:lagrangian_dynamics}).
This schedule avoids the instability of fixed-$\lambda$ training: too low permits unconstrained masks; too high collapses the predictor to trivial uniform scores.
The deadzone $\epsilon$ is essential---without it, the penalty oscillates around the target and interferes with KL minimization during the critical Gumbel annealing phase (Appendix~\ref{app:lagrangian}).

\subsection{Gradient Analysis: BCE vs.\ KL}
\label{sec:method:gradient}

\begin{figure}[t]
\centering
\includegraphics[width=\columnwidth]{snr_trajectory.pdf}
\vspace{-8pt}
\caption{Gradient SNR trajectory across 1000 training steps. BCE (orange) maintains consistently higher SNR than KL (blue) from step 100 onward ($2.3$--$17\times$), yet KL produces superior masks (Table~\ref{tab:main}), confirming that target alignment matters more than gradient clarity.}
\label{fig:snr_trajectory}
\vspace{-8pt}
\end{figure}

We contrast the gradient structure of BCE and KL to build intuition for how KL discovers divergent mask patterns (\S\ref{sec:mask_analysis}).

Under BCE, the predictor trains against oracle masks $m_j^* = \mathbf{1}[|g_j| \geq \text{top-}k\text{ threshold}]$ derived from activation magnitudes.
The gradient $\partial \mathcal{L}_{\text{BCE}} / \partial s_j = \sigma(s_j) - m_j^*$ depends only on neuron $j$; each neuron is trained independently with no inter-neuron coupling.
Under output-distribution KL, the chain rule introduces a downstream Jacobian that couples neuron $j$'s gradient to all other retained neurons in layers $\ell$ through $L$, evaluating each neuron in the context of the full sparse computation.

Gradient SNR measurements\footnote{$\text{SNR}_\ell = \|\mathbb{E}[\nabla_{\theta_\ell} \mathcal{L}]\|_2^2 / \text{Var}[\nabla_{\theta_\ell} \mathcal{L}]$, expectation over mini-batches.} reveal that BCE maintains consistently higher SNR than KL throughout training ($2.3\times$ at convergence; Figure~\ref{fig:snr_trajectory}; Appendix~\ref{app:snr_extended}), indicating KL's advantage stems not from gradient clarity but from discovering superior mask patterns via end-to-end search.
The target-swap ablation in \S\ref{sec:mask_analysis} confirms that the discovered mask targets, not the coupling mechanism itself, drive inference-stage quality.

\subsection{Two-Dimensional Taxonomy of Sparsity Prediction Methods}
\label{sec:method:taxonomy}

We organize prior work along two axes: \emph{correction granularity} (how the mask adapts to input) and \emph{training signal coupling} (whether the objective captures inter-neuron dependencies).

\begin{table}[t]
\centering
\caption{Taxonomy of learned sparsity prediction methods along correction granularity (rows) and training signal coupling (columns). Our method occupies the Per-token $\times$ Coupled cell (organizational framework). Empty cells represent unexplored combinations; their relative merit is an open question.}
\label{tab:taxonomy}
\resizebox{\columnwidth}{!}{
\begin{tabular}{lccc}
\toprule
 & Decomposed (BCE) & Intermediate (MSE) & Coupled (KL) \\
\midrule
Static & --- & --- & SPON, MaskLLM \\
Per-token & DejaVu, FastForward & ShadowLLM, D2DMoE & \textbf{Ours} \\
\bottomrule
\end{tabular}}
\end{table}

Our method combines per-token prediction with coupled KL training (Table~\ref{tab:taxonomy}).
KL on output logits directly optimizes next-token prediction loss, which composes across layers; per-layer MSE does not: a mask minimizing FFN output error at layer $\ell$ may amplify errors in layers $\ell{+}1 \ldots L$, explaining its $1.7\times$/$9.6\times$ worse PPL (Appendix~\ref{app:taxonomy}).

